% ============================================================================
% ENGLISH ABSTRACT
% ============================================================================
\begin{latin}
\begin{center}
    {\Large \textbf{Abstract}}
\end{center}
\vspace{0.5cm}

In the digital era, color images require specialized encryption algorithms due to their high data volume, multi-channel structure, and strong statistical correlation between adjacent pixels, making traditional text-based methods such as AES ineffective. Chaotic systems are suitable for image encryption due to their extreme sensitivity to initial conditions and ability to generate pseudo-random sequences; however, classical systems like Logistic and Tent maps suffer from "chaos degradation" in digital implementations, which reduces algorithm security. In this study, a five-stage algorithm for color image encryption is introduced, which leverages a combination of the three-dimensional Chen chaotic system and nine exponential chaotic maps (ECMs). In the first stage, the initial conditions and system parameters with a floating-point precision of 15 decimal digits are defined as the encryption key. In the second stage, the positions of image pixels are randomly shuffled (permutated) using the chaotic sequence obtained from the Chen system. In the third stage, one of the nine exponential chaotic maps is independently selected for each pixel. In the fourth stage, the color values of the pixels are combined with the generated key bytes. In the fifth stage, to break any remaining linear dependency and enhance resistance against differential and chosen-plaintext attacks, an additional diffusion layer with an independent Logistic map is applied, guaranteeing complete diffusion of changes across the image. The proposed algorithm was evaluated on a comprehensive set of standard images from the USC-SIPI and Kodak reference databases of various dimensions. Experimental results demonstrated that the average Shannon entropy of the encrypted images reached \lr{7.9973} bits, which is very close to the theoretical ideal value (8 bits) with a deviation of only \lr{0.0027} bits — consistent with natural statistical fluctuations in chaotic pseudo-random sequences. The correlation coefficient of adjacent pixels was reduced from an average of 0.9989 in the original images to 0.0082 in the encrypted images. Additionally, the number of pixel change rate (NPCR) and unified average changing intensity (UACI) achieved average values of 96.50\% and 30.93\%, respectively, and the complete reversibility of the algorithm was verified. Key space analysis indicates that the key space of the algorithm exceeds $2^{350}$ (equivalent to more than $10^{105}$), rendering any brute-force attack computationally impossible. The main innovation of this research is the dynamic and independent selection of the chaotic map for each pixel, which increases the structural complexity and security stability of the system.

\vspace{1cm}
\textbf{Keywords:} Image Encryption, Exponential Chaotic Systems, Chen Chaotic System, Permutation and Diffusion, Shannon Entropy, Key Space.
\end{latin}
\newpage
